\begin{frame}[label=q5b1]{5b.1 · DataFrame: promedio por departamento}
  \small
  \pregunta{¿Cuál es el promedio de personas contabilizadas por registro en cada departamento?}
  \entrada{93\,831 filas. \texttt{departamento} y \texttt{cantidad}; \texttt{count(*)} aporta el denominador.}
  \flujo{Registros por\\departamento}{\texttt{groupBy}\\Departamento}{\texttt{avg(cantidad)}\\Redondeo a 3 decimales}{26 departamentos\\Promedio descendente}
  \begin{columns}[T,onlytextwidth]
    \column{0.55\textwidth}
    \textbf{Código Scala original (fragmento)}
    \codigo{sql/05b_1_promedio_por_departamento}{12}{17}
    \column{0.42\textwidth}
    \textbf{Salida}\\[5pt]
    {\footnotesize
\begin{tabular}{@{}lll@{}}
\toprule
Departamento & Promedio & Personas\\
\midrule
Callao & 1.505 & 4\,565\\
Lima Met. & 1.485 & 39\,490\\
Lambayeque & 1.319 & 6\,231\\
\bottomrule
\end{tabular}
    }
  \end{columns}
  \respuesta{Callao encabeza el promedio, aunque Lima Metropolitana tiene mayor total. La media describe la cantidad por cruce agregado, no el tamaño de un incidente.}
  \origen{sql}{05b_1_promedio_por_departamento}
\end{frame}
